% ID: SEL-D0063-OB-E03
% ESTADO: cotejada. Segunda pasada visual y borrador compilado revisados.
% Fuente: Selectividad/sel_2010_matematicas_aplicadas/mat_apli/mat_apli_junio/mat_apli_junio_rb/EXAMEN-ANDALUCÍA-6.pdf
% SHA-256: 0de2f7d077a3b5d5c0dbca886cdb363082356411843369c3bce54ff7cda08abc
% Página: 2; opción B; ejercicio original 3.

Una fábrica posee un sistema de alarma contra robos. Por estudios previos a la instalación del sistema se sabe que la probabilidad de que un día se produzca un robo en la fábrica es $0.08$.

Las indicaciones técnicas del fabricante de la alarma dicen que la probabilidad de que suene si se ha producido un robo es $0.98$, y de que suene si no ha habido robo es $0.03$.
\begin{enumerate}[label=\alph*)]
\item \textbf{(1.25 puntos)} En un día cualquiera calcule la probabilidad de que no suene la alarma.
\item \textbf{(1.25 puntos)} Si suena la alarma, \textquestiondown{}cuál es la probabilidad de que no sea debido a un robo?
\end{enumerate}
